
\documentclass[12pt,a4paper]{report}

% Packages
\usepackage{amsmath}
\usepackage{amssymb}
\usepackage{graphicx}
\usepackage{braket}
\usepackage{booktabs}
\usepackage{array}
\usepackage{geometry}

% Page margins
\geometry{a4paper, margin=1in}

% Document information
\title{Monte Carlo Wave Function Simulation}
\author{}
\date{}

\begin{document}

\maketitle

\chapter{Details}

\section{Introduction to the MCWF Method}

The Monte Carlo Wave Function (MCWF) method is used to
study the time evolution of open quantum systems.

In this method, the quantum state evolves under a
non-Hermitian effective Hamiltonian between randomly
occurring quantum jumps.

The initial numerical implementation focuses on
non-Hermitian time evolution. It provides a foundation
for implementing the complete MCWF method, which also
includes stochastic quantum jumps and ensemble averaging.

\section{Numerical Simulation of a Two-Level Quantum System}

This simulation studies the time evolution of a driven
two-level quantum system in the presence of decay.

A non-Hermitian effective Hamiltonian is used to describe
the evolution between quantum jumps. The time-dependent
Schrödinger equation is solved numerically using the
Euler method.

The resulting state populations are plotted to study
the dynamics of the system.

The time-dependent Schrödinger equation is

\begin{equation}
i\hbar\frac{d}{dt}\ket{\psi(t)}
=
H_{\mathrm{eff}}\ket{\psi(t)}.
\end{equation}

Using natural units where $\hbar=1$,

\begin{equation}
\frac{d}{dt}\ket{\psi(t)}
=
-iH_{\mathrm{eff}}\ket{\psi(t)}.
\end{equation}

For a small time interval $\Delta t$, the Euler method
gives

\begin{equation}
\boxed{
\ket{\psi(t+\Delta t)}
\approx
\left(I-iH_{\mathrm{eff}}\Delta t\right)
\ket{\psi(t)}
}
\end{equation}

This equation is used to calculate the state at each
time step.

\section{Collapse Operator}

The collapse operator describes the effect of a
quantum jump caused by the interaction of the system
with its environment.

For a two-level atom, spontaneous emission causes
a transition from the excited state to the ground state:

\begin{equation}
\ket{e}\longrightarrow\ket{g}.
\end{equation}

The lowering operator responsible for this transition is

\begin{equation}
\sigma_-=\ket{g}\bra{e}.
\end{equation}

The collapse operator is defined as

\begin{equation}
\boxed{
C=\sqrt{\gamma}\,\sigma_-
}
\end{equation}

where $\gamma$ is the spontaneous-emission decay rate.

The factor $\sqrt{\gamma}$ introduces the decay rate
into the collapse operator.

The effective non-Hermitian Hamiltonian is

\begin{equation}
\boxed{
H_{\mathrm{eff}}
=
H-\frac{i}{2}C^\dagger C
}
\end{equation}

where $H$ is the Hermitian Hamiltonian of the system.

Since

\begin{equation}
C^\dagger C
=
\gamma\ket{e}\bra{e},
\end{equation}

the effective Hamiltonian becomes

\begin{equation}
\boxed{
H_{\mathrm{eff}}
=
H-\frac{i\gamma}{2}\ket{e}\bra{e}
}
\end{equation}

Here, we have used $\hbar=1$.

The non-Hermitian term causes the norm of the
unnormalized state to decrease during evolution
without a quantum jump.

\section{Quantum Jump}

A quantum jump is a sudden, stochastic change in the
quantum state of a system caused by its interaction
with the environment.

In the MCWF method, quantum jumps occur randomly
according to probabilities determined by the
instantaneous quantum state.

\subsection{Observation of Quantum Jumps}

In the simulated quantum trajectory, the excited-state
population shows sudden drops at approximately

\begin{equation}
\boxed{
t\approx
1.8,\ 3.1,\ 9.4,\ 11.7,\ 14.1,\ 17.7
}
\end{equation}

These sudden drops are consistent with spontaneous-emission
quantum jumps, provided that the simulation explicitly
implements stochastic collapse events.

For example, at approximately $t=17.7$, the excited-state
population decreases from about $0.8$ to $0$.

This corresponds to the transition

\begin{equation}
\ket{e}
\xrightarrow{\text{spontaneous emission}}
\ket{g}.
\end{equation}

After a spontaneous-emission jump, the quantum state
is reset to the ground state.

\subsection{Why Does the Excited-State Population Increase Again?}

After each quantum jump, the excited-state population
starts increasing again.

This happens because the ground and excited states
are coherently coupled by the driving field.

The coupling strength is represented by $\Omega_{12}$.

The driving field can excite the atom again:

\begin{equation}
\ket{g}
\xrightarrow{\Omega_{12}}
\ket{e}.
\end{equation}

The excited atom may then undergo spontaneous emission:

\begin{equation}
\ket{e}
\xrightarrow{\gamma}
\ket{g}.
\end{equation}

This process of coherent excitation and spontaneous
emission can repeat throughout the simulation.

\subsection{Why Do Quantum Jumps Occur at Random Times?}

In the MCWF method, the probability of a quantum jump
during a small time interval $dt$ depends on the
instantaneous quantum state.

For a normalized state

\begin{equation}
\ket{\psi(t)}
=
c_g(t)\ket{g}+c_e(t)\ket{e},
\end{equation}

the probability of a spontaneous-emission jump is

\begin{equation}
P_{\mathrm{jump}}
=
dt\bra{\psi(t)}C^\dagger C\ket{\psi(t)}.
\end{equation}

Substituting

\begin{equation}
C=\sqrt{\gamma}\ket{g}\bra{e},
\end{equation}

we obtain

\begin{equation}
\boxed{
P_{\mathrm{jump}}
=
\gamma |c_e(t)|^2dt
}
\end{equation}

to first order in $dt$.

Here,

\begin{itemize}
    \item $\gamma$ is the spontaneous-emission decay rate.
    \item $|c_e(t)|^2$ is the excited-state population.
    \item $dt$ is the small time interval.
\end{itemize}

The occurrence of a quantum jump is determined
probabilistically rather than at fixed time intervals.

Therefore, jumps generally occur at irregular times.

A quantum jump does not require the excited-state
population to reach unity. It may occur whenever
the excited-state population is nonzero.

\section{Ensemble Averaging of Quantum Trajectories}

A single MCWF trajectory contains randomly occurring
quantum jumps. Therefore, its population dynamics
can show sudden changes.

To obtain the average behaviour of the system,
multiple independent trajectories are simulated.

The ensemble-averaged population is calculated as

\begin{equation}
\boxed{
\overline{P_e(t)}
=
\frac{1}{N}
\sum_{k=1}^{N}P_e^{(k)}(t)
}
\end{equation}

where $N$ is the total number of trajectories and
$P_e^{(k)}(t)$ is the excited-state population of
the $k$-th trajectory.

Similarly, the ground-state population is

\begin{equation}
\boxed{
\overline{P_g(t)}
=
\frac{1}{N}
\sum_{k=1}^{N}P_g^{(k)}(t)
}
\end{equation}

For normalized two-level states,

\begin{equation}
\overline{P_g(t)}+\overline{P_e(t)}=1.
\end{equation}

\section{Results and Discussion}

The population dynamics are compared for different
numbers of MCWF trajectories.

For a small number of trajectories, the averaged
population curves show noticeable statistical
fluctuations because the quantum jumps occur randomly.

As the number of trajectories increases from 10 to
10,000, these fluctuations become smaller, and
the averaged population curves become smoother.

In the simulation, the ground-state and excited-state
populations approach approximately

\begin{equation}
\boxed{
P_g\approx0.67,\qquad P_e\approx0.33
}
\end{equation}

at long times.

These values indicate the approximate steady-state
populations for the chosen simulation parameters.

Increasing the number of trajectories improves the
statistical accuracy of the ensemble average.

For independent trajectories, the statistical error
typically decreases approximately as

\begin{equation}
\boxed{
\text{Statistical error}\propto\frac{1}{\sqrt{N}}
}
\end{equation}

where $N$ is the number of trajectories.

\section{Conclusion}

The numerical simulation demonstrates the evolution
of a driven two-level quantum system in the presence
of spontaneous emission.

The effective non-Hermitian Hamiltonian describes
the evolution between quantum jumps, while the
collapse operator describes spontaneous-emission events.

The random occurrence of quantum jumps makes each
individual trajectory different.

However, averaging over a large number of trajectories
reduces statistical fluctuations and produces smoother
population curves.

The results illustrate the main features of the
Monte Carlo Wave Function method and show how
ensemble averaging can be used to study the
dynamics of an open quantum system.


\section{Future}

In this project, I implemented the Monte Carlo Wave Function method to study the dynamics of a driven two-level quantum system. The simulation captures random quantum jumps due to spontaneous emission and the coherent evolution between the ground and excited states. By increasing the number of trajectories from 10 to 10,000, I observed that the statistical fluctuations decreased and the populations approached steady values.

As future work, I plan to extend this simulation to specific single-qubit gates, such as X, Y, Z, and Hadamard, and study how environmental noise affects their fidelity. I also plan to compare the results with the Lindblad master equation for numerical validation.

% pdflatex -interaction=nonstopmode -halt-on-error -output-directory=build details.tex
\end{document}
